\documentclass[11pt]{article}
\usepackage[margin=1in]{geometry}
\usepackage{amsmath,amssymb,amsthm}
\usepackage[T1]{fontenc}
\usepackage{lmodern}
\usepackage{microtype}
\usepackage[colorlinks=true,linkcolor=blue,citecolor=blue,urlcolor=blue]{hyperref}
\newcommand{\Z}{\mathbb Z}
\newcommand{\C}{2^{\mathbb N}}
\newcommand{\Act}{\operatorname{Act}}
\newcommand{\asdimB}{\operatorname{asdim}_{\mathrm B}}
\newtheorem{theorem}{Theorem}
\newtheorem{lemma}[theorem]{Lemma}
\hypersetup{pdftitle={Generic Cantor actions of Z wr Z are hyperfinite},pdfauthor={Alec Kriebel}}
\title{Generic Cantor actions of $\Z\wr\Z$ are hyperfinite}
\author{Alec Kriebel\\\small Independent researcher\\\small\href{https://orcid.org/0009-0001-9320-500X}{ORCID: 0009-0001-9320-500X}}
\date{October 13, 2026\quad (version 1.0.0)}
\begin{document}
\maketitle
\begin{abstract}
We give an affirmative answer to Question 3 in Iyer's contribution, joint with
Shinko, to Oberwolfach Report 2/2025. A dense $G_\delta$ of continuous actions
of the restricted wreath product $\Z\wr\Z$ on Cantor space has hyperfinite
orbit equivalence relation on the entire space. We specialize established
Cantor-action lifting principles to obtain a continuous profinite shift height.
Finite-width lamp slabs then satisfy the compatible-metric theorem of
Conley--Jackson--Marks--Seward--Tucker-Drob. The full preimage of the embedded
integers is treated by a separate finite Borel exhaustion. The auxiliary Borel
theorem requires the height factor and allows stabilizers.
\end{abstract}

\section{The question and the imported theorems}
Let $N=\bigoplus_{i\in\Z}\Z$ and $W=N\rtimes\Z$, with lamp generators $a_i$
and shift $t$ satisfying $ta_it^{-1}=a_{i+1}$. Put $a=a_0$ and let
$\chi:W\to\Z$ be the shift homomorphism. Here $\wr$ denotes the restricted,
countable discrete wreath product. Give $\Act(W,\C)$ the subspace topology of
$\operatorname{Homeo}(\C)^W$, with compact-open homeomorphism topologies.
It is a closed Polish space. Generic means comeager in this full action space.
A countable Borel equivalence relation is \emph{hyperfinite} if it is the
increasing union of finite Borel equivalence relations.

\begin{theorem}\label{main}
There is a dense $G_\delta$ subset of $\Act(W,\C)$ such that every action in
it has hyperfinite orbit equivalence relation on all of $\C$.
\end{theorem}

This is the exact whole-space question posed as open in
\cite[printed p.~115, Question~3]{OWR}. The locally finite-asymptotic-dimension
theorem of Iyer--Shinko \cite{IS} does not cover $W$: it is finitely generated
and contains undistorted copies of $\Z^d$ for arbitrarily large $d$.
Our theorem makes no restriction to a conull or comeager subset of the phase
space, and does not assert hyperfiniteness for every Borel $W$-action.

We use three published consequences of Conley--Jackson--Marks--Seward--
Tucker-Drob \cite{CJMST}, with numbering from their author manuscript.
Corollary~5.5 gives finite Borel asymptotic dimension for \emph{any} Borel
action of each fixed $\Z^m$. Lemmas~2.2 and~3.1 and the paragraph following
Definition~3.2 give its finite-color characterization: for every finite set of
group moves there is a Borel coloring with a fixed finite number of colors
and uniformly finite monochromatic move components. The number of colors
depends on $m$, not on the move set. Finite Borel covers may be disjointized
to obtain colorings.

Theorem~7.3 of \cite{CJMST} says that if proper Borel extended metrics
$\rho_n$ on a standard Borel space each have finite Borel asymptotic
dimension and
\begin{equation}\label{compatibility}
\sup\{\rho_{n+1}(x,y):\rho_n(x,y)<R\}<\infty
\quad(n\geq1,\ R>0),
\end{equation}
then $\bigcup_n E_{\rho_n}$ is hyperfinite, where $E_\rho$ is the relation of
finite $\rho$-distance. No common bound on the dimensions is required.
Finally, Corollary~7.5 gives hyperfiniteness for \emph{any} Borel action of
a countable group of local polynomial growth; it applies to $N$.
These are deep imported results, rather than claims proved by our diagnostics.

\section{Generic compatible heights}
Write $\widehat\Z=\varprojlim_n\Z/n!\Z$, with $W$ acting by addition of
$\chi$. The generic-height argument below is an explicit specialization of
established lifting principles. Akin--Glasner--Weiss \cite[Theorem~3.3,
equation~(3.7)]{AGW} use conditional compatible lifts for single
homeomorphisms. Doucha \cite[Proposition~2.12(3)]{Doucha} supplies all-group
extension perturbation. Doucha--Melleray--Tsankov
\cite[Proposition~3.2, Definition~A.3, Proposition~A.7]{DMT} give generic
finite-extension realization for all Cantor group actions. In their Boolean
coding, the finite partition equations below are a specialization of that
principle. We include the direct construction to check the prescribed
compatibility and the full action-space topology.

\begin{lemma}\label{height}
A dense $G_\delta$ of $\Act(W,\C)$ admits a continuous equivariant map
$p:\C\to\widehat\Z$.
\end{lemma}
\begin{proof}
Put $\mathcal A=\Act(W,\C)$. For continuous $c:\C\to\Z/m\Z$, let
$D(c,m)$ be the set of actions $\alpha$ for which
\[
c(\alpha(a)x)=c(x),\qquad c(\alpha(t)x)=c(x)+1\pmod m.
\]
These equations imply $c(\alpha(g)x)=c(x)+\chi(g)$ for every $g\in W$.
The set $D(c,m)$ is clopen: the finite clopen fibers of $c$ give uniform
neighborhoods preserving all output labels of both generator maps.
There are countably many such $c$, since every clopen set is a finite union
from a countable clopen basis. For $m\mid n$, define the open set
\[
U(c;m,n)=(\mathcal A\setminus D(c,m))\ \cup\!
\bigcup_{d\bmod m=c}D(d,n),
\]
where the union ranges over continuous maps $d:\C\to\Z/n\Z$.

To prove density, fix $\alpha\in D(c,m)$ and an action neighborhood $V$.
The fiber product
\[
Y=\{(x,j)\in\C\times\Z/n\Z:c(x)=j\bmod m\}
\]
carries the exact continuous action
$\gamma(g)(x,j)=(\alpha(g)x,j+\chi(g))$. Choose a finite symmetric set $F$
of group elements controlling $V$, a finite clopen partition $Q$ of
sufficiently small mesh, and a common refinement $P$ of $Q$, all
$\alpha(g)^{-1}Q$ for $g\in F$, and the fibers of $c$.
For each nonempty atom $A$ of $P$, both $A$ and the inverse image $Y_A$
of $A$ are Cantor spaces: $Y_A$ consists of $n/m$ copies of $A$.
Choose atomwise homeomorphisms $A\to Y_A$ and combine them into
$h:\C\to Y$. Let $\beta(g)=h^{-1}\gamma(g)h$.

If $h(x)=(u,j)$, then $u,x$ are in the same $P$-atom. For $g\in F$,
$\alpha(g)u$ and $\alpha(g)x$ are in one $Q$-atom. Also
$\beta(g)x$ and $\alpha(g)u$ lie in one $P$-atom, hence one $Q$-atom.
Thus $\beta(g)$ is uniformly close to $\alpha(g)$ on all of $F$.
Symmetry of $F$ controls inverse maps as well, so $\beta\in V$.
The second coordinate $d$ of $h$ is a continuous equivariant residue map
modulo $n$, and $d\bmod m=c$ because $P$ refines the old fibers.
This proves density of $U(c;m,n)$; outside $D(c,m)$ membership is automatic.

Intersect these open dense sets over all $c,m,n$ with $m\mid n$.
For a fixed action in the resulting dense $G_\delta$, start with $c_1=0$
and recursively lift $c_{n!}$ to $c_{(n+1)!}$. Every map selected in this
recursion was already included in the countably many obligations.
The compatible residues define continuous equivariant $p$.
Its compact nonempty image is invariant under $\pm1$, whose integer
translates are dense in $\widehat\Z$, so $p$ is onto. A measurable choice
of $p$ as the action varies is not required.
\end{proof}

\section{The Borel height theorem}
\begin{theorem}\label{borel}
Suppose $W$ acts Borelly on a standard Borel space $X$ and there is a Borel
map $p:X\to\widehat\Z$ with $p(gx)=p(x)+\chi(g)$. Then its entire orbit
equivalence relation $E_W$ is hyperfinite.
\end{theorem}

No freeness, minimality, invariant probability, or surjectivity of $p$ is
assumed in Theorem~\ref{borel}.

\subsection*{Finite-width lamp slabs}
For $m\geq1$, let $r_m(x)\in\{0,\ldots,m-1\}$ be the residue of $p(x)$ and
$B_m=r_m^{-1}(0)$. This is an $N$-invariant Borel set. The coordinate map
\[
x\longleftrightarrow (z,r)=(t^{-r_m(x)}x,r_m(x))\in B_m\times\{0,\ldots,m-1\}
\]
is a Borel bijection with inverse $(z,r)\mapsto t^r z$. Uniqueness uses
the prescribed residue, even when lamps have stabilizers.

Let $G_m$ be the undirected graph including all lamp edges $\{x,ax\}$ and
the shift edge $\{x,tx\}$ when $r_m(x)<m-1$, with loops omitted.
Its degree is at most four. Each distance-at-most-$k$ relation is a finite
union of graphs of partial Borel move words, so its path metric $\rho_m$
is a proper Borel extended metric. Proper here means that finite-radius
balls are finite. In the displayed coordinates its edges are
\[
(z,r)\longleftrightarrow(a_{-r}z,r),\qquad
(z,r)\longleftrightarrow(z,r+1)\quad(r<m-1),
\]
since $t^{-r}at^r=a_{-r}$. The subgroup
$K_m=\langle a_0,a_{-1},\ldots,a_{-(m-1)}\rangle\cong\Z^m$ acts on $B_m$.

\begin{lemma}\label{slab}
Each $\rho_m$ has finite Borel asymptotic dimension.
\end{lemma}
\begin{proof}
By \cite[Corollary~5.5]{CJMST} and its finite-move characterization, there
is a fixed finite color count $D_m+1$ such that, for each $R\geq1$, $B_m$
has a Borel coloring whose monochromatic components under $K_m$ words
of length at most $\lceil R\rceil$ have at most $C_R<\infty$ points.
Lift the color of $z$ to every $(z,r)$.
A slab path of length at most $R$ changes its base by at most
$\lceil R\rceil$ lamp letters. Hence a monochromatic $\rho_m$-distance-at-most-$R$
chain projects into one monochromatic base component. Connecting paths
need not themselves be monochromatic; only their endpoints are used.
The lifted component has at most $mC_R$ points, and its diameter is at
most $R(mC_R-1)$ by a simple path in the threshold graph.
For $0<R<1$ the components are singletons. This is the finite-color
characterization of finite Borel asymptotic dimension in
\cite[Lemma~3.1]{CJMST}. No bound uniform in $m$ is asserted.
\end{proof}

\subsection*{Compatible exhaustion on the regular part}
If $m\mid M$, a forbidden residue $M-1$ reduces to the forbidden residue
$m-1$. Thus $G_m\subseteq G_M$ and $\rho_M\leq\rho_m$ on all pairs of
finite old distance. In particular factorial widths satisfy
\eqref{compatibility}, with bound $R$.

The embedded integers form a countable Borel subset of $\widehat\Z$.
Put $X_{\rm exc}=p^{-1}(\Z)$ and $X_{\rm reg}=X\setminus X_{\rm exc}$;
both are Borel and $W$-invariant. On $X_{\rm reg}$, every forward shift
edge eventually appears in a factorial slab: a point with
$r_{n!}(x)=n!-1$ for every $n$ would have $p(x)=-1$, contrary to membership
in this set. Once admitted, an edge stays admitted. Lamp edges are always
present. Each finite generator word therefore fits in one sufficiently
late slab, giving
\[
E_W\mathbin{\upharpoonright}X_{\rm reg}
=\bigcup_n E_{\rho_{n!}\mathbin{\upharpoonright}X_{\rm reg}}.
\]
Invariance means that restricting the graphs removes no intervening path.
Lemma~\ref{slab} and \cite[Theorem~7.3]{CJMST} imply hyperfiniteness here.
This invokes the compatible-metric theorem, not a general increasing-union
principle for hyperfinite relations.

\subsection*{The full exceptional preimage}
For $x\in X_{\rm exc}$, let $h(x)$ be the unique integer equal to $p(x)$,
and put $B=p^{-1}(0)$ and $q(x)=t^{-h(x)}x\in B$.
Both $h$ and $q$ are Borel, and $N$ preserves $B$.
For $y=gx$, the element $t^{-h(y)}gt^{h(x)}$ has zero shift character and
lies in $N$. Conversely an $N$-move between normalized points conjugates
back to a $W$-move. Therefore
\[
x\ E_W\ y\quad\Longleftrightarrow\quad q(x)\ E_N\ q(y).
\]
Every finitely generated subgroup of $N$ has polynomial growth, so
\cite[Corollary~7.5]{CJMST} supplies increasing finite Borel relations
$F_j$ on $B$ exhausting its $N$ orbit relation. On $X_{\rm exc}$ define
\[
x\ H_j\ y\quad\Longleftrightarrow\quad
x=y\ \text{ or }\
\bigl(|h(x)|,|h(y)|\leq j\ \text{ and }\ q(x)\ F_j\ q(y)\bigr).
\]
Inside the band this is a pullback of $F_j$; outside it all classes are
singletons. For each $z\in B$ and $k\in\Z$ there is exactly one point
with $q(x)=z$ and $h(x)=k$, namely $t^kz$. Thus band classes are finite,
of size at most $(2j+1)|[z]_{F_j}|$. The relations $H_j$ are Borel,
increase, and exhaust the full exceptional orbit relation.
The exceptional set is not assumed countable or negligible; pulling back
$F_j$ without the height band would not give finite classes.

Take the disjoint union of finite increasing exhaustions of the two
invariant pieces. This proves Theorem~\ref{borel} on all of $X$.
Lemma~\ref{height} then proves Theorem~\ref{main}.

\section{Verification and publication status}
The accompanying package preserves the original proof and reviews, adds
fresh independent mathematical and dated priority audits, and supplies
portable finite diagnostics and integrity checks. Finite computations
check group identities, nonfree slab models, divisor cuts, projections,
and bounded-height exhaustions; they do not prove the infinite theorem.
The proof above relies on the precise published inputs stated in Section~1.

The bounded priority audit of October~13, 2026 found no exact earlier
resolution in its inspected corpus and identified the lifting predecessors
credited above. It does not certify historical firstness or worldwide
current openness. This note answers the historically posed exact question.

OpenAI tools were used extensively in developing and checking the argument,
drafting the note, writing diagnostics, and conducting independent adversarial
AI reviews. This is an unrefereed preprint. No completed external human peer
review or formal proof-assistant verification is claimed.

\begin{thebibliography}{9}
\bibitem{AGW} E.~Akin, E.~Glasner, and B.~Weiss,
\emph{Generically there is but one self homeomorphism of the Cantor set},
Trans. Amer. Math. Soc. \textbf{360} (2008), 3613--3630.
\href{https://doi.org/10.1090/S0002-9947-08-04450-4}{doi:10.1090/S0002-9947-08-04450-4}.
\bibitem{CJMST} C.~T.~Conley, S.~C.~Jackson, A.~S.~Marks, B.~M.~Seward,
and R.~D.~Tucker-Drob,
\emph{Borel asymptotic dimension and hyperfinite equivalence relations},
Duke Math. J. \textbf{172} (2023), 3175--3226.
\href{https://doi.org/10.1215/00127094-2022-0100}{doi:10.1215/00127094-2022-0100}.
Author manuscript: \url{https://math.berkeley.edu/~marks/papers/polycyclic_v13.pdf}.
\bibitem{Doucha} M.~Doucha,
\emph{Strong topological Rokhlin property, shadowing, and symbolic dynamics
of countable groups}, \href{https://arxiv.org/abs/2211.08145v3}{arXiv:2211.08145v3} (2024).
\bibitem{DMT} M.~Doucha, J.~Melleray, and T.~Tsankov,
\emph{Dense and comeager conjugacy classes in zero-dimensional dynamics},
\href{https://arxiv.org/abs/2507.05474v1}{arXiv:2507.05474v1} (2025);
revised \href{https://arxiv.org/abs/2507.05474v2}{v2} (2026).
\bibitem{OWR} S.~Iyer (joint with F.~Shinko),
\emph{Generic actions of countable groups}, in \emph{Set Theory},
Oberwolfach Report 2/2025, 114--115, Question~3.
\href{https://doi.org/10.4171/owr/2025/2}{doi:10.4171/owr/2025/2}.
\bibitem{IS} S.~Iyer and F.~Shinko,
\emph{Asymptotic dimension and hyperfiniteness of generic Cantor actions},
Groups Geom. Dyn. (2025), published online first.
\href{https://doi.org/10.4171/GGD/896}{doi:10.4171/GGD/896}.
Inspected manuscript: \href{https://arxiv.org/abs/2409.03078v1}{arXiv:2409.03078v1}.
\end{thebibliography}
\end{document}
